% The argument as a closed LOOP: build the torus into a cubic along the top
% (T -> wp,wp' -> E), then a line on E hands three points back to the torus whose
% sum is zero (bottom), and the return arrow IS the group law.  Teal = the three
% spaces; gold = the two mechanisms; the rust return arrow is the punchline.
% Drawn at native scale; poster.tex sizes it to the block width.  (shadings loaded;
% no shadows library -- depth comes from a subtle top/bottom gradient.)
\begin{tikzpicture}[
    >=Latex, line cap=round, line join=round,
    obj/.style={rectangle, rounded corners=7pt, draw=formteal!72!black, line width=1.2pt,
      top color=formteal!3, bottom color=formteal!20, align=center,
      inner sep=9pt, minimum width=5.8cm, minimum height=2.6cm},
    fact/.style={rectangle, rounded corners=7pt, draw=formgold!85!black, line width=1.2pt,
      top color=formgold!4, bottom color=formgold!26, align=center,
      inner sep=9pt, minimum width=5.8cm, minimum height=2.6cm},
    flow/.style={-{Latex[length=4mm,width=3mm]}, line width=1.7pt, draw=formslate},
    back/.style={-{Latex[length=4.6mm,width=3.4mm]}, line width=2.2pt, draw=formrust!88!black},
    chip/.style={fill=white, draw=black!16, rounded corners=2.5pt, inner sep=4pt,
      font=\small, align=center},
    achip/.style={fill=formrust!9, draw=formrust!55, rounded corners=2.5pt, inner sep=5pt,
      font=\small\bfseries, text=formrust!72!black, align=center},
  ]
  % --- the three spaces (top row) ---
  \node[obj] (T)  at (0,0)  {\textbf{complex torus}\\[2pt]$T=\C/\Lambda$\\[1pt]{\footnotesize\itshape addition is free}};
  \node[obj] (wp) at (9,0)  {\textbf{Weierstrass}\\\textbf{functions} $\wp,\wp'$};
  \node[obj] (E)  at (18,0) {\textbf{plane cubic}\\[2pt]$E\subset\CP^2$\\[1pt]{\footnotesize$y^2=4x^3-g_2x-g_3$}};
  % --- the two mechanisms (bottom row) ---
  \node[fact] (line) at (18,-7.4) {\textbf{a line meets }$\boldsymbol{E}$\\[2pt]in three points $P,Q,R$};
  \node[fact] (div)  at (0,-7.4)  {\textbf{their torus-sum is }$\boldsymbol{0}$\\[2pt]$P\oplus Q\oplus R=O$};
  % --- central tie: a focal badge linking the two rows ---
  \node[rounded corners=6pt, draw=formteal!45!black, line width=1pt, align=center,
        fill=white, fill opacity=0.9, text opacity=1, inner sep=8pt, text=formteal!42!black]
        at (9,-3.85)
    {\itshape the loop closes:\\[3pt]{\large$\C/\Lambda\;\cong\;E$}\\[2pt]\itshape= the group law};
  % --- arrows around the loop ---
  \draw[flow] (T)  -- node[chip,above=1pt] {$\Phi$\\[1pt]{\footnotesize$[\wp{:}\wp'{:}1]$}} (wp);
  \draw[flow] (wp) -- node[chip,above=1pt] {lands\\on $E$} (E);
  \draw[flow] (E)  -- node[chip,left=3pt] {cut by\\a line} (line);
  \draw[flow] (line) -- node[chip,below=1pt] {Abel relation:\ \ $\sum z_i\equiv 0 \ (\mathrm{mod}\ \Lambda)$} (div);
  \draw[back] (div) -- node[achip,right=3pt] {collinear\\$\Leftrightarrow$ sum $=O$} (T);
\end{tikzpicture}
